% GENERATED FILE. Edit canonical lessons, then run npm run generate:books.
% canonical-source-hash: fnv1a64:84bf97c1a93511fa
% canonical-lessons: JA-C322-shouchi-suru, JA-C322-ukeireru, JA-C322-mochiron, JA-C322-zehi, JA-C322-moushideru, JA-R322-review-317-322

\chapter{Saying Yes}
\label{ch:ja-saying-yes}
\hlchaptermodality{322}

\begin{chapteropening}
\textbf{By the end of this chapter:} \emph{I can agree to a request and say why.}

Complete the last lesson of chapter 322: I can agree to a request and say why.
\end{chapteropening}

\section[\texorpdfstring{\ja{しょうちする} (shōchi suru)}{しょうちする (shōchi suru)}]{\ja{しょうちする} (shōchi suru) — to consent; to agree to a request}

\label{lesson:JA-C322-shouchi-suru}



\begin{quote}
Before the new one: say the Japanese for rather, then the Japanese for to nod.
\end{quote}

\subsection*{You'll want to know: \ja{しょうちする}}
\textbf{\ja{しょうちする}} — \emph{shōchi suru} — \textquotedblleft{}to consent; to agree to a request\textquotedblright{}.

\begin{quote}
\textbf{\ja{あにが} \ja{たのんだので、しょうちしました。}} — \emph{ani ga tanonda node, shōchi shimashita} — \textquotedblleft{}My brother asked me, so I agreed.\textquotedblright{}
\end{quote}

Polite \textbf{\ja{しょうちします}}, \emph{shōchishimasu}; a \ja{する} verb. Past (polite): \textbf{\ja{てんいんが} \ja{しょうちしました。}}, \emph{ten'in ga shōchi shimashita}.

\begin{cousinweb}
\textbf{\ja{しょうちする}} (\emph{shōchi suru}) is built from a Sino-Japanese noun and the native verb \emph{suru}, \textquotedblleft{}to do\textquotedblright{}. The noun, \emph{shōchi}, is two Chinese characters read with their on'yomi, readings borrowed from Chinese: \emph{shō}, \textquotedblleft{}acquiesce\textquotedblright{} (Middle Chinese \emph{dzying}) and \emph{chi}, \textquotedblleft{}know\textquotedblright{} (\emph{trje}) (KANJIDIC2; JMdict; Middle Chinese in Baxter's transcription).
\end{cousinweb}

\begin{grammarlens}[title={What you've built}]
One more word for giving reasons and opinions.
\end{grammarlens}

\subsection*{Your turn}
\begin{itemize}
\raggedright
  \item \emph{Say it:} \emph{shōchi suru}
  \item \emph{Say it:} \emph{shōchi suru}, once more
  \item \emph{Say it:} \emph{ani ga tanonda node, shōchi shimashita} — My brother asked me, so I agreed.
  \item \emph{Recall:} say the Japanese for rather, then the Japanese for to nod, then say \emph{shōchi suru} again
\end{itemize}

\begin{tcolorbox}[breakable,colback=teal!4,colframe=teal!35!black,title={Before you move on}]
What does \ja{しょうちする} mean? (\textquotedblleft{}To consent; to agree to a request\textquotedblright{}.) Say it once more.
\end{tcolorbox}
\section[\texorpdfstring{\ja{うけいれる} (ukeireru)}{うけいれる (ukeireru)}]{\ja{うけいれる} (ukeireru) — to accept; to take in}

\label{lesson:JA-C322-ukeireru}



\begin{quote}
Before the new one: say the Japanese for to nod, then the Japanese for to consent.
\end{quote}

\subsection*{You'll want to know: \ja{うけいれる}}
\textbf{\ja{うけいれる}} — \emph{ukeireru} — \textquotedblleft{}to accept; to take in\textquotedblright{}.

\begin{quote}
\textbf{\ja{みんなが} \ja{さんせいしたので、かいしゃは} \ja{ていあんを} \ja{うけいれました。}} — \emph{minna ga sansei shita node, kaisha wa teian o ukeiremashita} — \textquotedblleft{}Everyone agreed, so the company accepted the proposal.\textquotedblright{}
\end{quote}

Polite \textbf{\ja{うけいれます}}, \emph{ukeiremasu}; an ichidan verb. Past (polite): \textbf{\ja{かいしゃが} \ja{うけいれました。}}, \emph{kaisha ga ukeiremashita}.

\begin{cousinweb}
\textbf{\ja{うけいれる}} (\emph{ukeireru}) is a native Japanese word (wago), inherited rather than borrowed from Chinese: it is built from \emph{ukeru}, \textquotedblleft{}to receive\textquotedblright{} and \emph{ireru}, \textquotedblleft{}to put in\textquotedblright{}, and its kanji were attached later as kun'yomi (JMdict).
\end{cousinweb}

\begin{grammarlens}[title={What you've built}]
One more word for giving reasons and opinions.
\end{grammarlens}

\subsection*{Your turn}
\begin{itemize}
\raggedright
  \item \emph{Say it:} \emph{ukeireru}
  \item \emph{Say it:} \emph{ukeireru}, once more
  \item \emph{Say it:} \emph{minna ga sansei shita node, kaisha wa teian o ukeiremashita} — Everyone agreed, so the company accepted the proposal.
  \item \emph{Recall:} say the Japanese for to nod, then the Japanese for to consent, then say \emph{ukeireru} again
\end{itemize}

\begin{tcolorbox}[breakable,colback=teal!4,colframe=teal!35!black,title={Before you move on}]
What does \ja{うけいれる} mean? (\textquotedblleft{}To accept; to take in\textquotedblright{}.) Say it once more.
\end{tcolorbox}
\section[\texorpdfstring{\ja{もちろん} (mochiron)}{もちろん (mochiron)}]{\ja{もちろん} (mochiron) — of course}

\label{lesson:JA-C322-mochiron}



\begin{quote}
Before the new one: say the Japanese for to consent, then the Japanese for to take in.
\end{quote}

\subsection*{You'll want to know: \ja{もちろん}}
\textbf{\ja{もちろん}} — \emph{mochiron} — \textquotedblleft{}of course\textquotedblright{}.

\begin{quote}
\textbf{\ja{じかんが} \ja{あるから、もちろん} \ja{てつだいます。}} — \emph{jikan ga aru kara, mochiron tetsudaimasu} — \textquotedblleft{}I have time, so of course I will help.\textquotedblright{}
\end{quote}

\begin{cousinweb}
\textbf{\ja{もちろん}} (\emph{mochiron}) is Sino-Japanese (kango): two Chinese characters read with their on'yomi, readings borrowed from Chinese. They are \emph{mochi}, \textquotedblleft{}not\textquotedblright{} and \emph{ron}, \textquotedblleft{}argument\textquotedblright{} (Middle Chinese \emph{lwon}) (KANJIDIC2; Middle Chinese in Baxter's transcription). JMdict notes that it is usually written in kana.
\end{cousinweb}

\begin{grammarlens}[title={What you've built}]
One more word for giving reasons and opinions.
\end{grammarlens}

\subsection*{Your turn}
\begin{itemize}
\raggedright
  \item \emph{Say it:} \emph{mochiron}
  \item \emph{Say it:} \emph{mochiron}, once more
  \item \emph{Say it:} \emph{jikan ga aru kara, mochiron tetsudaimasu} — I have time, so of course I will help.
  \item \emph{Recall:} say the Japanese for to consent, then the Japanese for to take in, then say \emph{mochiron} again
\end{itemize}

\begin{tcolorbox}[breakable,colback=teal!4,colframe=teal!35!black,title={Before you move on}]
What does \ja{もちろん} mean? (\textquotedblleft{}Of course\textquotedblright{}.) Say it once more.
\end{tcolorbox}
\section[\texorpdfstring{\ja{ぜひ} (zehi)}{ぜひ (zehi)}]{\ja{ぜひ} (zehi) — by all means; without fail}

\label{lesson:JA-C322-zehi}



\begin{quote}
Before the new one: say the Japanese for to take in, then the Japanese for of course.
\end{quote}

\subsection*{You'll want to know: \ja{ぜひ}}
\textbf{\ja{ぜひ}} — \emph{zehi} — \textquotedblleft{}by all means; without fail\textquotedblright{}.

\begin{quote}
\textbf{\ja{おもしろいので、ぜひ} \ja{よんで} \ja{ください。}} — \emph{omoshiroi node, zehi yonde kudasai} — \textquotedblleft{}It is interesting, so do read it.\textquotedblright{}
\end{quote}

\begin{cousinweb}
\textbf{\ja{ぜひ}} (\emph{zehi}) is Sino-Japanese (kango): two Chinese characters read with their on'yomi, readings borrowed from Chinese. They are \emph{ze}, \textquotedblleft{}right\textquotedblright{} (Middle Chinese \emph{dzyeX}) and \emph{hi}, \textquotedblleft{}negative\textquotedblright{} (KANJIDIC2; Middle Chinese in Baxter's transcription).
\end{cousinweb}

\begin{grammarlens}[title={What you've built}]
One more word for giving reasons and opinions.
\end{grammarlens}

\subsection*{Your turn}
\begin{itemize}
\raggedright
  \item \emph{Say it:} \emph{zehi}
  \item \emph{Say it:} \emph{zehi}, once more
  \item \emph{Say it:} \emph{omoshiroi node, zehi yonde kudasai} — It is interesting, so do read it.
  \item \emph{Recall:} say the Japanese for to take in, then the Japanese for of course, then say \emph{zehi} again
\end{itemize}

\begin{tcolorbox}[breakable,colback=teal!4,colframe=teal!35!black,title={Before you move on}]
What does \ja{ぜひ} mean? (\textquotedblleft{}By all means; without fail\textquotedblright{}.) Say it once more.
\end{tcolorbox}
\section[\texorpdfstring{\ja{もうしでる} (mōshideru)}{もうしでる (mōshideru)}]{\ja{もうしでる} (mōshideru) — to come forward; to offer}

\label{lesson:JA-C322-moushideru}



\begin{quote}
Before the new one: say the Japanese for of course, then the Japanese for by all means.
\end{quote}

\subsection*{You'll want to know: \ja{もうしでる}}
\textbf{\ja{もうしでる}} — \emph{mōshideru} — \textquotedblleft{}to come forward; to offer\textquotedblright{}.

\begin{quote}
\textbf{\ja{あには} \ja{じかんが} \ja{あったので、もうしでました。}} — \emph{ani wa jikan ga atta node, mōshidemashita} — \textquotedblleft{}My brother had time, so he offered to help.\textquotedblright{}
\end{quote}

Polite \textbf{\ja{もうしでます}}, \emph{mōshidemasu}; an ichidan verb. Past (polite): \textbf{\ja{あにが} \ja{もうしでました。}}, \emph{ani ga mōshidemashita}.

\begin{cousinweb}
\textbf{\ja{もうしでる}} (\emph{mōshideru}) is a native Japanese word (wago), inherited rather than borrowed from Chinese: it is built from \emph{mōsu}, \textquotedblleft{}to say (humble)\textquotedblright{} and \emph{deru}, \textquotedblleft{}to come out\textquotedblright{}, and its kanji were attached later as kun'yomi (JMdict).
\end{cousinweb}

\begin{grammarlens}[title={What you've built}]
One more word for giving reasons and opinions.
\end{grammarlens}

\subsection*{Your turn}
\begin{itemize}
\raggedright
  \item \emph{Say it:} \emph{mōshideru}
  \item \emph{Say it:} \emph{mōshideru}, once more
  \item \emph{Say it:} \emph{ani wa jikan ga atta node, mōshidemashita} — My brother had time, so he offered to help.
  \item \emph{Recall:} say the Japanese for of course, then the Japanese for by all means, then say \emph{mōshideru} again
\end{itemize}

\begin{tcolorbox}[breakable,colback=teal!4,colframe=teal!35!black,title={Before you move on}]
What does \ja{もうしでる} mean? (\textquotedblleft{}To come forward; to offer\textquotedblright{}.) Say it once more.
\end{tcolorbox}
\section[(dialogue)]{Looking back: from Agreeing and Disagreeing to Saying Yes}

\label{lesson:JA-R322-review-317-322}



\begin{quote}
Say the Japanese for to agree and to be against.
\end{quote}

\subsection*{Your turn}
Say these aloud:

\begin{itemize}
\raggedright
  \item to agree, to be against, a discussion, to propose, a conclusion
  \item to explain, to be convinced, grounds for a claim, a meaning, for what reason
  \item one's real feelings, the official line, to admit, to misunderstand, in earnest
  \item a point of view, a doubt, to insist, to comprehend, especially
  \item absolutely, without fail, as expected, rather, to nod
  \item to consent, to take in, of course, by all means, to come forward
\end{itemize}

\begin{tcolorbox}[breakable,colback=teal!4,colframe=teal!35!black,title={Before you move on}]
To agree? (\textbf{\ja{さんせいする}}, \emph{sansei suru}.) To be against? (\textbf{\ja{はんたいする}}, \emph{hantai suru}.) A discussion? (\textbf{\ja{ぎろん}}, \emph{giron}.) To propose? (\textbf{\ja{ていあんする}}, \emph{teian suru}.) A conclusion? (\textbf{\ja{けつろん}}, \emph{ketsuron}.) To explain? (\textbf{\ja{せつめいする}}, \emph{setsumei suru}.) To be convinced? (\textbf{\ja{なっとくする}}, \emph{nattoku suru}.) Grounds for a claim? (\textbf{\ja{こんきょ}}, \emph{konkyo}.) A meaning? (\textbf{\ja{いみ}}, \emph{imi}.) For what reason? (\textbf{\ja{なぜ}}, \emph{naze}.) One's real feelings? (\textbf{\ja{ほんね}}, \emph{honne}.) The official line? (\textbf{\ja{たてまえ}}, \emph{tatemae}.) To admit? (\textbf{\ja{みとめる}}, \emph{mitomeru}.) To misunderstand? (\textbf{\ja{ごかいする}}, \emph{gokai suru}.) In earnest? (\textbf{\ja{ほんき}}, \emph{honki}.) A point of view? (\textbf{\ja{たちば}}, \emph{tachiba}.) A doubt? (\textbf{\ja{ぎもん}}, \emph{gimon}.) To insist? (\textbf{\ja{しゅちょうする}}, \emph{shuchō suru}.) To comprehend? (\textbf{\ja{りかいする}}, \emph{rikai suru}.) Especially? (\textbf{\ja{とくに}}, \emph{tokuni}.) Absolutely? (\textbf{\ja{ぜったい}}, \emph{zettai}.) Without fail? (\textbf{\ja{かならず}}, \emph{kanarazu}.) As expected? (\textbf{\ja{やはり}}, \emph{yahari}.) Rather? (\textbf{\ja{むしろ}}, \emph{mushiro}.) To nod? (\textbf{\ja{うなずく}}, \emph{unazuku}.) To consent? (\textbf{\ja{しょうちする}}, \emph{shōchi suru}.) To take in? (\textbf{\ja{うけいれる}}, \emph{ukeireru}.) Of course? (\textbf{\ja{もちろん}}, \emph{mochiron}.) By all means? (\textbf{\ja{ぜひ}}, \emph{zehi}.) To come forward? (\textbf{\ja{もうしでる}}, \emph{mōshideru}.)
\end{tcolorbox}
